\chapter{The Passive Voice}

The active voice focuses on the person or thing performing an action. The
passive focuses on the action or its result. The performer can be omitted when
it is unknown, obvious, or unimportant.

\begin{center}
\textit{Die Mechanikerin repariert das Auto.}
\quad $\longrightarrow$ \quad
\textit{Das Auto wird repariert.}
\end{center}

\section{Process passive: \textit{werden + Partizip II}}

The process passive (\textit{Vorgangspassiv}) describes an action or process.
The finite form of \textit{werden} carries the tense, and the Partizip II
stands at the end.

\begin{center}
\small
\rowcolors{2}{referenceBlueBg}{white}
\begin{tabularx}{\textwidth}{@{}>{\bfseries}l l X@{}}
\toprule
Tense & Formation & Example \\
\midrule
Präsens & wird + Partizip II & Das Auto \textbf{wird repariert}. \\
Präteritum & wurde + Partizip II & Das Auto \textbf{wurde repariert}. \\
Perfekt & ist + Partizip II + worden & Das Auto \textbf{ist repariert worden}. \\
Plusquamperfekt & war + Partizip II + worden & Das Auto \textbf{war repariert worden}. \\
Futur I & wird + Partizip II + werden & Das Auto \textbf{wird repariert werden}. \\
Modal verb & modal + Partizip II + werden & Das Auto \textbf{muss repariert werden}. \\
\bottomrule
\end{tabularx}
\end{center}

In the passive Perfekt, use \textbf{worden}, not \textit{geworden}:
\textit{Das Haus ist gebaut worden}. By contrast, \textit{geworden} is the
Partizip II of \textit{werden} meaning ``to become'':
\textit{Er ist müde geworden}.

\section{What happens to the objects?}

An accusative object in the active sentence becomes the nominative subject of
the passive sentence. A dative object remains dative.

\begin{center}
\small
\rowcolors{2}{referenceGreenBg}{white}
\begin{tabularx}{\textwidth}{@{}>{\bfseries}l X@{}}
\toprule
Active & Passive \\
\midrule
Man repariert \textbf{das Auto}.
  & \textbf{Das Auto} wird repariert. \\
Man gibt \textbf{dem Gast} \textbf{den Schlüssel}.
  & \textbf{Der Schlüssel} wird \textbf{dem Gast} gegeben. \\
Man hilft \textbf{dem Mann}.
  & \textbf{Dem Mann} wird geholfen. \\
\bottomrule
\end{tabularx}
\end{center}

The final example has no nominative subject because \textit{helfen} governs
the dative. This subjectless passive is normal German.

\section{Naming the performer or cause}

Use \textit{von + Dativ} for a person, institution, or originator. Use
\textit{durch + Akkusativ} especially for a means, cause, or process.

\begin{itemize}[leftmargin=*]
  \item \textit{Der Brief wurde \textbf{von der Chefin} geschrieben.}
  \item \textit{Die Straße wurde \textbf{durch den Sturm} blockiert.}
\end{itemize}

\section{State passive: \textit{sein + Partizip II}}

The state passive (\textit{Zustandspassiv}) describes the condition resulting
from an action, not the action itself.

\begin{center}
\small
\rowcolors{2}{referenceYellowBg}{white}
\begin{tabularx}{\textwidth}{@{}>{\bfseries}l l X@{}}
\toprule
Meaning & Formation & Example \\
\midrule
Process & werden + Partizip II & Die Tür \textbf{wird geschlossen}. \\
Resulting state & sein + Partizip II & Die Tür \textbf{ist geschlossen}. \\
Past process & wurde + Partizip II & Die Tür \textbf{wurde geschlossen}. \\
Past state & war + Partizip II & Die Tür \textbf{war geschlossen}. \\
\bottomrule
\end{tabularx}
\end{center}

\section{Impersonal passive}

An activity can be presented without naming any performer. Initial
\textit{es} disappears when another element occupies position one.

\begin{center}
\textit{Es wird heute getanzt.}
\qquad
\textit{Heute wird getanzt.}

\textit{Im Unterricht wird viel gesprochen.}
\end{center}

\section{Alternatives to the passive}

German often uses \textit{man} when the performer is general, or
\textit{sich lassen + infinitive} to express possibility.

\begin{itemize}[leftmargin=*]
  \item Passive: \textit{Das Problem kann gelöst werden.}
  \item With \textit{man}: \textit{Man kann das Problem lösen.}
  \item With \textit{sich lassen}: \textit{Das Problem lässt sich lösen.}
\end{itemize}

\begin{tcolorbox}[
  colback=referenceRedBg,
  colframe=genderFeminine,
  title={Word-order reminder},
  fonttitle=\bfseries
]
In a subordinate clause, the finite passive auxiliary moves to the end:
\textit{Ich weiß, dass das Auto heute \textbf{repariert wird}.}
With a modal verb: \textit{Ich weiß, dass das Auto heute
\textbf{repariert werden muss}.}
\end{tcolorbox}
